% GENERATED FILE. Edit canonical lessons, then run npm run generate:books.
% canonical-source-hash: fnv1a64:82314f8f8a86e42f
% canonical-lessons: KA-C284-hunnime, KA-C284-suryasta, KA-C284-ninne, KA-C284-samaya, KA-C284-kelavomme

\chapter{Full moon, Sunset, Yesterday, Time, Sometimes}
\label{ch:ka-full-moon-sunset-yesterday-time-sometimes}
\hlchaptermodality{284}

\begin{chapteropening}
\textbf{By the end of this chapter:} \emph{I can say full moon, sunset, yesterday, time and sometimes.}

Complete the last lesson of chapter 284: I can say full moon, sunset, yesterday, time and sometimes.
\end{chapteropening}

\section[\texorpdfstring{\kn{ಹುಣ್ಣಿಮೆ} (huṇṇime)}{ಹುಣ್ಣಿಮೆ (huṇṇime)}]{\kn{ಹುಣ್ಣಿಮೆ} (huṇṇime) — full moon (day)}

\label{lesson:KA-C284-hunnime}



\begin{quote}
Before the new one: say the Kannada for an operation, then the Kannada for pregnant.
\end{quote}

\subsection*{You'll want to know: \kn{ಹುಣ್ಣಿಮೆ}}
\textbf{\kn{ಹುಣ್ಣಿಮೆ}} — \emph{huṇṇime} — \textquotedblleft{}full moon (day)\textquotedblright{}.

\begin{grammarlens}[title={What you've built}]
One more word for this chapter.
\end{grammarlens}

\subsection*{Your turn}
\begin{itemize}
\raggedright
  \item \emph{Say it:} \emph{huṇṇime}
  \item \emph{Say it:} \emph{huṇṇime}, once more
  \item \emph{Say it:} say \emph{huṇṇime}
  \item \emph{Recall:} say the Kannada for an operation, then the Kannada for pregnant, then say \emph{huṇṇime} again
\end{itemize}

\begin{tcolorbox}[breakable,colback=teal!4,colframe=teal!35!black,title={Before you move on}]
What does \kn{ಹುಣ್ಣಿಮೆ} mean? (\textquotedblleft{}Full moon (day)\textquotedblright{}.) Say it once more.
\end{tcolorbox}
\section[\texorpdfstring{\kn{ಸೂರ್ಯಾಸ್ತ} (sūryāsta)}{ಸೂರ್ಯಾಸ್ತ (sūryāsta)}]{\kn{ಸೂರ್ಯಾಸ್ತ} (sūryāsta) — sunset}

\label{lesson:KA-C284-suryasta}



\begin{quote}
Before the new one: say the Kannada for pregnant, then the Kannada for full moon.
\end{quote}

\subsection*{You'll want to know: \kn{ಸೂರ್ಯಾಸ್ತ}}
\textbf{\kn{ಸೂರ್ಯಾಸ್ತ}} — \emph{sūryāsta} — \textquotedblleft{}sunset\textquotedblright{}.

\begin{grammarlens}[title={What you've built}]
One more word for this chapter.
\end{grammarlens}

\subsection*{Your turn}
\begin{itemize}
\raggedright
  \item \emph{Say it:} \emph{sūryāsta}
  \item \emph{Say it:} \emph{sūryāsta}, once more
  \item \emph{Say it:} say \emph{sūryāsta}
  \item \emph{Recall:} say the Kannada for pregnant, then the Kannada for full moon, then say \emph{sūryāsta} again
\end{itemize}

\begin{tcolorbox}[breakable,colback=teal!4,colframe=teal!35!black,title={Before you move on}]
What does \kn{ಸೂರ್ಯಾಸ್ತ} mean? (\textquotedblleft{}Sunset\textquotedblright{}.) Say it once more.
\end{tcolorbox}
\section[\texorpdfstring{\kn{ನಿನ್ನೆ} (ninne)}{ನಿನ್ನೆ (ninne)}]{\kn{ನಿನ್ನೆ} (ninne) — yesterday}

\label{lesson:KA-C284-ninne}



\begin{quote}
Before the new one: say the Kannada for full moon, then the Kannada for sunset.
\end{quote}

\subsection*{You'll want to know: \kn{ನಿನ್ನೆ}}
\textbf{\kn{ನಿನ್ನೆ}} — \emph{ninne} — \textquotedblleft{}yesterday\textquotedblright{}.

\begin{grammarlens}[title={What you've built}]
One more word for this chapter.
\end{grammarlens}

\subsection*{Your turn}
\begin{itemize}
\raggedright
  \item \emph{Say it:} \emph{ninne}
  \item \emph{Say it:} \emph{ninne}, once more
  \item \emph{Say it:} say \emph{ninne}
  \item \emph{Recall:} say the Kannada for full moon, then the Kannada for sunset, then say \emph{ninne} again
\end{itemize}

\begin{tcolorbox}[breakable,colback=teal!4,colframe=teal!35!black,title={Before you move on}]
What does \kn{ನಿನ್ನೆ} mean? (\textquotedblleft{}Yesterday\textquotedblright{}.) Say it once more.
\end{tcolorbox}
\section[\texorpdfstring{\kn{ಸಮಯ} (samaya)}{ಸಮಯ (samaya)}]{\kn{ಸಮಯ} (samaya) — time}

\label{lesson:KA-C284-samaya}



\begin{quote}
Before the new one: say the Kannada for sunset, then the Kannada for yesterday.
\end{quote}

\subsection*{You'll want to know: \kn{ಸಮಯ}}
\textbf{\kn{ಸಮಯ}} — \emph{samaya} — \textquotedblleft{}time\textquotedblright{}.

\begin{grammarlens}[title={What you've built}]
One more word for this chapter.
\end{grammarlens}

\subsection*{Your turn}
\begin{itemize}
\raggedright
  \item \emph{Say it:} \emph{samaya}
  \item \emph{Say it:} \emph{samaya}, once more
  \item \emph{Say it:} say \emph{samaya}
  \item \emph{Recall:} say the Kannada for sunset, then the Kannada for yesterday, then say \emph{samaya} again
\end{itemize}

\begin{tcolorbox}[breakable,colback=teal!4,colframe=teal!35!black,title={Before you move on}]
What does \kn{ಸಮಯ} mean? (\textquotedblleft{}Time\textquotedblright{}.) Say it once more.
\end{tcolorbox}
\section[\texorpdfstring{\kn{ಕೆಲವೊಮ್ಮೆ} (kelavomme)}{ಕೆಲವೊಮ್ಮೆ (kelavomme)}]{\kn{ಕೆಲವೊಮ್ಮೆ} (kelavomme) — sometimes}

\label{lesson:KA-C284-kelavomme}



\begin{quote}
Before the new one: say the Kannada for yesterday, then the Kannada for time.
\end{quote}

\subsection*{You'll want to know: \kn{ಕೆಲವೊಮ್ಮೆ}}
\textbf{\kn{ಕೆಲವೊಮ್ಮೆ}} — \emph{kelavomme} — \textquotedblleft{}sometimes\textquotedblright{}.

\begin{grammarlens}[title={What you've built}]
One more word for this chapter.
\end{grammarlens}

\subsection*{Your turn}
\begin{itemize}
\raggedright
  \item \emph{Say it:} \emph{kelavomme}
  \item \emph{Say it:} \emph{kelavomme}, once more
  \item \emph{Say it:} say \emph{kelavomme}
  \item \emph{Recall:} say the Kannada for yesterday, then the Kannada for time, then say \emph{kelavomme} again
\end{itemize}

\begin{tcolorbox}[breakable,colback=teal!4,colframe=teal!35!black,title={Before you move on}]
What does \kn{ಕೆಲವೊಮ್ಮೆ} mean? (\textquotedblleft{}Sometimes\textquotedblright{}.) Say it once more.
\end{tcolorbox}
